\documentclass[11pt,oneside,openany]{book}

\usepackage[T1]{fontenc}
\usepackage[utf8]{inputenc}
\usepackage{newtxtext}
\usepackage{amsmath}
\usepackage{amsthm}
\usepackage{newtxmath}
\usepackage[scaled=0.92]{inconsolata}
\usepackage{microtype}
\usepackage[letterpaper,margin=1in,headheight=15pt]{geometry}
\usepackage{mathtools,bm}
\usepackage{booktabs,longtable,tabularx,array,multirow}
\usepackage{enumitem}
\usepackage{xcolor}
\usepackage{listings}
\usepackage{tikz}
\usetikzlibrary{arrows.meta,positioning,calc,fit,backgrounds,shapes.geometric,decorations.pathreplacing,matrix}
\usepackage[most]{tcolorbox}
\usepackage{float}
\usepackage{algorithm}
\usepackage{algpseudocode}
\makeatletter
\renewcommand{\theALG@line}{\thealgorithm.\arabic{ALG@line}}
\renewcommand{\theHALG@line}{\thealgorithm.\arabic{ALG@line}}
\makeatother
\usepackage{fancyhdr}
\usepackage{caption}
\usepackage{subcaption}
\usepackage{hyperref}
\usepackage[nameinlink,noabbrev]{cleveref}

\definecolor{bsblue}{HTML}{24557A}
\definecolor{bslightblue}{HTML}{EAF2F8}
\definecolor{bsgreen}{HTML}{3A7D44}
\definecolor{bslightgreen}{HTML}{EDF7EE}
\definecolor{bsorange}{HTML}{B5651D}
\definecolor{bslightorange}{HTML}{FFF3E6}
\definecolor{bsred}{HTML}{A33A3A}
\definecolor{bsgray}{HTML}{5E6670}
\definecolor{bslightgray}{HTML}{F3F5F7}
\definecolor{bsdark}{HTML}{1F2933}
\definecolor{cbBlue}{HTML}{0072B2}
\definecolor{cbOrange}{HTML}{E69F00}
\definecolor{cbPurple}{HTML}{CC79A7}

\hypersetup{
  colorlinks=true,
  linkcolor=bsblue,
  citecolor=bsgreen,
  urlcolor=bsorange,
  pdftitle={EllipticBSAMTree Code Documentation},
  pdfauthor={Documentation of the EllipticBSAMTree and CHStokesBSAMTree MATLAB implementations},
  pdfsubject={Strict-tree block-structured adaptive multigrid: elliptic and Cahn-Hilliard-Stokes solvers}
}

\pagestyle{fancy}
\fancyhf{}
\fancyhead[L]{\small EllipticBSAMTree Code Documentation}
\fancyhead[R]{\small\nouppercase{\leftmark}}
\fancyfoot[C]{\thepage}
\renewcommand{\headrulewidth}{0.4pt}

\setcounter{secnumdepth}{3}
\setcounter{tocdepth}{2}
\setlength{\parindent}{0pt}
\setlength{\parskip}{5pt plus 1pt minus 1pt}
\emergencystretch=2em

\newcommand{\code}[1]{\textnormal{\texttt{#1}}}
\newcommand{\R}{\mathbb{R}}
\newcommand{\dd}{\,\mathrm{d}}
\DeclareMathOperator{\divop}{div}
\newcommand{\ULap}{\mathrm{ULap}}
\newcommand{\UDiv}{\mathrm{UDiv}}

\newtheorem{invariant}{Implementation invariant}[chapter]
\newtheorem{remark}{Remark}[chapter]

\newtcolorbox{keyidea}[1][]{
  enhanced,breakable,colback=bslightblue,colframe=bsblue,
  boxrule=0.7pt,arc=1.5mm,left=2mm,right=2mm,top=1.5mm,bottom=1.5mm,
  title={Key idea},fonttitle=\bfseries,#1}
\newtcolorbox{codepath}[1][]{
  enhanced,breakable,colback=bslightgreen,colframe=bsgreen,
  boxrule=0.7pt,arc=1.5mm,left=2mm,right=2mm,top=1.5mm,bottom=1.5mm,
  title={Code path},fonttitle=\bfseries,#1}
\newtcolorbox{cautionbox}[1][]{
  enhanced,breakable,colback=bslightorange,colframe=bsorange,
  boxrule=0.7pt,arc=1.5mm,left=2mm,right=2mm,top=1.5mm,bottom=1.5mm,
  title={Implementation note},fonttitle=\bfseries,#1}
\newtcolorbox{fortranbox}[1][]{
  enhanced,breakable,colback=bslightgray,colframe=bsgray,
  boxrule=0.7pt,arc=1.5mm,left=2mm,right=2mm,top=1.5mm,bottom=1.5mm,
  title={Fortran correspondence},fonttitle=\bfseries,#1}

\lstdefinestyle{matlabstyle}{
  language=Matlab,
  basicstyle=\small\ttfamily,
  keywordstyle=\color{bsblue}\bfseries,
  commentstyle=\color{bsgreen!70!black},
  stringstyle=\color{bsorange!90!black},
  numbers=left,
  numberstyle=\scriptsize\color{bsgray},
  stepnumber=1,
  numbersep=8pt,
  frame=single,
  rulecolor=\color{bsgray!35},
  backgroundcolor=\color{bslightgray},
  breaklines=true,
  breakatwhitespace=false,
  showstringspaces=false,
  columns=fullflexible,
  keepspaces=true,
  tabsize=2,
  captionpos=b
}
\lstset{style=matlabstyle}

\title{\bfseries EllipticBSAMTree Code Documentation\\[6pt]
\Large A strict-tree block-structured adaptive multigrid solver\\
for cell-centered elliptic problems}
\author{Documentation of the supplied MATLAB implementation\\
\small following the Fortran BSAM 2.0 library (SuperBSAM)}
\date{July 2026}

\begin{document}
\frontmatter
\maketitle
\tableofcontents

\chapter*{Preface}
\addcontentsline{toc}{chapter}{Preface}

\code{EllipticBSAMTree} (function prefix \code{ebt}) solves
$-\divop(D\nabla u)+Cu=f$ on a rectangle with Dirichlet/Neumann sides,
using the block-structured adaptive FAS multigrid (BSAM) methodology of
Feng--Guo--Lowengrub--Wise \cite{feng2018} reorganized around a
\emph{strict single-parent patch tree} --- the data structure of the
Fortran reference library \code{SuperBSAM/BSAM2.0}
(\code{treeops.f90} and friends).  It is numerically identical to the
companion \code{EllipticBSAM} solver --- same discretization, same
QIF/LIF coarse--fine interpolation, same conservative flux corrections
--- but every fine patch has exactly one parent that fully contains it,
which turns every coarse--fine transfer into a single slice of one
parent array.

The reader is assumed to know the companion document \emph{BSAM Code
Documentation} (in the \code{EllipticBSAM} folder), which develops the
discretization, the QIF/LIF ghost interpolation, the flux-balance
corrections and the AFAS V-cycle in detail.  This document concentrates
on what is \emph{different}: the tree, its two index systems, per-parent
clustering, and the single-parent transfer machinery.  The two sibling
solvers built on this framework are documented in their own folders:
\code{CHStokesBSAMTree} (Cahn--Hilliard--Stokes) and \code{CHNSBSAMTree}
(Cahn--Hilliard--Navier--Stokes).

\mainmatter

% ======================================================================
\chapter{Orientation}
\label{ch:orientation}

\section{The solver family and its lineage}

Both packages inherit their numerics from \code{EllipticBSAM} and their
\emph{architecture} from the Fortran BSAM 2.0 library.  The family tree:

\begin{center}
\begin{tikzpicture}[
  box/.style={draw=bsblue, fill=bslightblue, rounded corners=1.5mm,
              minimum width=4.2cm, minimum height=1.05cm, align=center,
              font=\small},
  fbox/.style={draw=bsgray, fill=bslightgray, rounded corners=1.5mm,
              minimum width=4.2cm, minimum height=1.05cm, align=center,
              font=\small},
  arr/.style={-{Stealth[length=2.6mm]}, thick, bsdark}]
\node[box] (eb) {\code{EllipticBSAM}\\[-1pt]{\scriptsize QC-ring elliptic solver}};
\node[fbox, right=1.6cm of eb] (f) {\code{SuperBSAM/BSAM2.0}\\[-1pt]{\scriptsize Fortran reference (treeops)}};
\node[box, below=0.9cm of eb] (ebt) {\code{EllipticBSAMTree}\\[-1pt]{\scriptsize strict-tree elliptic solver}};
\node[box, below=0.9cm of f] (chs) {\code{CHStokesBSAMTree}\\[-1pt]{\scriptsize Cahn--Hilliard--Stokes port}};
\draw[arr] (eb) -- node[left, font=\scriptsize, align=right]{numerics\\(kernels, QIF)} (ebt);
\draw[arr] (f) -- node[right, pos=0.62, font=\scriptsize, align=left]{tree design\\(per-parent clustering)} (ebt);
\draw[arr] (ebt) -- node[below, font=\scriptsize, sloped]{tree framework} (chs);
\draw[arr] (f) -- node[right, font=\scriptsize, align=left]{operators, smoothers,\\time loop, output} (chs);
\node[box, right=1.5cm of chs] (chns) {\code{CHNSBSAMTree}\\[-1pt]{\scriptsize Navier--Stokes extension}};
\draw[arr] (chs) -- node[above, font=\scriptsize]{+ inertia, $\nu(\phi)$, $\rho(\phi)$} (chns);
\end{tikzpicture}
\end{center}

\section{What ``strict tree'' means}
\label{sec:stricttree}

In \code{EllipticBSAM} the clusterer works on a whole level at once, so
a fine patch may overlap several coarse patches; the per-patch
\code{parents\{p\}} field is a \emph{list}, and every coarse--fine
transfer gathers from a union of coarse patches through precomputed
per-cell index lists (the ``QC ring'' machinery).  That design is
correct and, as the companion audit showed, carries no convergence
penalty --- but it is not a tree.

The tree variants instead enforce, exactly as the Fortran does:

\begin{invariant}[single parent]
\label{inv:parent}
Every non-root patch is created by clustering \emph{inside one parent
patch}.  Its footprint, recorded in parent-local cell coordinates as
\code{mbounds}, satisfies $1 \le \code{mbounds} \le
\code{parent.mx}$: a child may touch its parent's edge but never extend
past it, and never spans two parents.
\end{invariant}

A physical feature that crosses two coarse patches therefore produces
\emph{two abutting children}, one per parent, which communicate as
same-level siblings --- never one straddling child
(\cref{fig:split}).

\begin{figure}[H]
\centering
\begin{tikzpicture}[scale=0.62]
% left: QC-ring behavior
\begin{scope}
\draw[thick, bsgray] (0,0) rectangle (4,4);
\draw[thick, bsgray] (4,0) rectangle (8,4);
\node[font=\scriptsize\color{bsgray}] at (2,3.6) {coarse patch $A$};
\node[font=\scriptsize\color{bsgray}] at (6,3.6) {coarse patch $B$};
\fill[bslightorange, draw=bsorange, thick] (2.5,1.2) rectangle (5.5,2.8);
\node[font=\scriptsize\color{bsorange}, align=center] at (4,2)
  {one child,\\two parents};
\node[font=\small\bfseries] at (4,-0.7) {\code{EllipticBSAM} (QC ring)};
\end{scope}
% right: tree behavior
\begin{scope}[xshift=10cm]
\draw[thick, bsgray] (0,0) rectangle (4,4);
\draw[thick, bsgray] (4,0) rectangle (8,4);
\fill[bslightblue, draw=bsblue, thick] (2.5,1.2) rectangle (4,2.8);
\fill[bslightgreen, draw=bsgreen, thick] (4,1.2) rectangle (5.5,2.8);
\node[font=\scriptsize\color{bsblue}, anchor=east] at (2.35,2) {child of $A$};
\node[font=\scriptsize\color{bsgreen}, anchor=west] at (5.65,2) {child of $B$};
\draw[thick, bsdark, densely dashed] (4,1.2) -- (4,2.8);
\node[font=\scriptsize] at (4,0.75) {sibling seam};
\node[font=\small\bfseries] at (4,-0.7) {tree variants (per-parent clustering)};
\end{scope}
\end{tikzpicture}
\caption{A tagged feature crossing two coarse patches.  The QC-ring
solver creates one multi-parent child; the tree solvers split it per
parent, and the two children couple through the same-level (sibling)
exchange.}
\label{fig:split}
\end{figure}

\begin{keyidea}
The payoff of \cref{inv:parent} is architectural: the $(c_{m_x}{+}2)
\times (c_{m_y}{+}2)$ parent window of a child --- interior
\emph{and} ghost ring --- becomes \emph{one contiguous slice of the
parent's padded array}.  Every coarse--fine transfer (window gather,
ring refresh, restriction, FAS loading, correction, coefficient
fill-down) is a sliced copy against a single known array.  The per-cell
multi-patch resolver machinery of the QC-ring design disappears from
the hot path.
\end{keyidea}

\section{Where things live}

\begin{center}\small
\begin{tabular}{@{}lll@{}}
\toprule
 & \code{EllipticBSAMTree} & \code{CHStokesBSAMTree} \\
\midrule
entry point & \code{ebtSolve} & \code{chsRun} \\
problem definition & \code{ebtProblem} & \code{chsProblem} \\
options & \code{ebtOptions} & \code{chsOptions} \\
test suite & \code{tests/runAllTreeTests} (T1--T9) & \code{tests/runChsTests} (C1--C6) \\
tree validator & \code{ebtTreeCheck} & \code{chsTreeCheck} \\
\bottomrule
\end{tabular}
\end{center}

Both run headless as
\code{matlab -batch "cd('<pkg>/tests'); <suite>"}.  Diagnostic messages
use the tags \code{[ebtsam]} and \code{[chs]}.

% ======================================================================
\chapter{The Strict Patch Tree}
\label{ch:tree}

\section{Two index systems per patch}
\label{sec:indexsystems}

Following the Fortran \code{nodeinfo} type, every patch carries two
coordinate records (\cref{fig:coords}):

\begin{itemize}[leftmargin=2em]
\item \code{lev.box(p,:)} $= [i_{lo}\, i_{hi}\, j_{lo}\, j_{hi}]$ ---
the patch's cell range in the \emph{global index space of its own
level} (the Fortran \code{mglobal}).  It drives everything
\emph{horizontal}: owner maps, sibling ghost exchange, tag buffering,
plotting.
\item \code{lev.mbounds(p,:)} $= [I_l\, I_h\, J_l\, J_h]$ --- the
patch's coarsened footprint in \emph{its parent's local cell space},
1-based (the Fortran \code{mbounds}).  It drives everything
\emph{vertical}: window/ring slices, restriction and load scatters,
coefficient fill-down, correction windows.
\item \code{lev.parent(p)} --- the index of \emph{the} parent patch on
level $k{-}1$ ($0$ on the coarsest level); \code{lev.children\{q\}}
lists the direct children.
\end{itemize}

The two systems are locked together: with $q=\code{parent}(p)$ and
$\code{qb}=\code{box}(q,:)$,
\[
\code{mbounds}(p,:) =
\Bigl[\tfrac{i_{lo}+1}{2},\ \tfrac{i_{hi}}{2},\
      \tfrac{j_{lo}+1}{2},\ \tfrac{j_{hi}}{2}\Bigr]
- \bigl[\code{qb}_1{-}1,\ \code{qb}_1{-}1,\ \code{qb}_3{-}1,\
        \code{qb}_3{-}1\bigr],
\]
which requires the global box to be odd-start/even-end --- automatic,
because a child is the exact $2\times$ refinement of a cluster box.
\code{ebtTreeCheck}/\code{chsTreeCheck} verify this identity, the
containment of \cref{inv:parent}, pairwise-disjoint siblings, the
children/covered bookkeeping, and the ring nesting, on every hierarchy
the test suites build.

\begin{figure}[H]
\centering
\begin{tikzpicture}[scale=0.55]
% parent patch grid
\draw[step=1, bsgray!40, thin] (0,0) grid (12,8);
\draw[very thick, bsgray] (0,0) rectangle (12,8);
\node[font=\scriptsize\color{bsgray}, anchor=south west] at (0.1,8.1)
  {parent patch $q$: local cells $1..12 \times 1..8$};
% child footprint
\fill[bslightblue, fill opacity=0.75] (3,2) rectangle (9,6);
\draw[very thick, bsblue] (3,2) rectangle (9,6);
\node[font=\scriptsize\color{bsblue}, align=center] at (6,4)
  {child footprint\\$\code{mbounds}=[4\ 9\ 3\ 6]$};
% ring
\draw[very thick, bsorange, densely dashed] (2,1) rectangle (10,7);
\node[font=\scriptsize\color{bsorange}, anchor=west] at (10.15,6.6)
  {window ring};
% brace for global coords
\draw[decorate, decoration={brace, mirror, amplitude=4pt}, thick]
  (3,-0.25) -- (9,-0.25)
  node[midway, below=5pt, font=\scriptsize]
  {child \code{box} (own level) $=[2\cdot 4{-}1+2(\code{qb}_1{-}1),\ \ldots]$};
\end{tikzpicture}
\caption{The lineage record.  The child's coarsened footprint
(\code{mbounds}, blue) lies inside its single parent; the window ring
(orange) extends one parent cell around it and never leaves the
parent's \emph{padded} array --- where it exits the parent's interior
it lands on the parent's own ghost cells.}
\label{fig:coords}
\end{figure}

\section{The single-parent window}
\label{sec:window}

Let the child have footprint $[I_l..I_h]\times[J_l..J_h]$ in parent
cells, $c_{m_x}=I_h{-}I_l{+}1$, $c_{m_y}=J_h{-}J_l{+}1$.  Since parent
cell $(i,j)$ sits at padded index $(i{+}1, j{+}1)$ and $1\le I_l$,
$I_h\le \code{parent.mx}$, the whole window
\[
W \;=\; \text{parent }Q\bigl(I_l : I_h{+}2,\ J_l : J_h{+}2\bigr)
\]
is a legal slice: rows $I_l$ and $I_h{+}2$ are the ring cells
$I_l{-}1$ and $I_h{+}1$, which are parent \emph{ghost} cells exactly
when the child touches the parent's edge.  This is the Fortran
\code{InterpolateCoarseFine} contract verbatim: \emph{the child reads
only \code{parent\%q}}; sibling and physical closures reach the child
indirectly, through the parent's synchronized ghost layer.

\begin{codepath}
\code{ebtBuildOps} stores, per patch: \code{ops.win} (the window
slice), \code{ops.ring} (its four perimeter slices, refreshed into
\code{QC} before each QIF interpolation), \code{ops.usc}/\code{ops.fsc}
(single-rect scatters of the restricted solution and the FAS load into
the parent), \code{ops.pcc/pdew/pdns} (parent coefficient slices), and
\code{ops.corrW/E/S/N} (flux-correction targets).  The V-cycle hot path
(\code{ebtTransfer}, \code{ebtFillGhosts}, \code{ebtCoarseLoad},
\code{ebtCoeffFilldown}) consists of sliced assignments against these
descriptors.
\end{codepath}

\section{Per-parent Berger--Rigoutsos clustering}
\label{sec:perparent}

Regridding follows the Fortran pipeline
(\code{bsamroutines.f90}: \code{AMR} $\to$ \code{EstimateLevelErrors}
$\to$ \code{GridAdapt} $\to$ \code{RefineGrid} $\to$
\code{NewSubGrids}) with one deliberate simplification
(\cref{sec:nesting}):

\begin{enumerate}[leftmargin=2em]
\item \textbf{Tag} on the current finest level, per patch
(indicator of choice; the \code{CHStokesBSAMTree} documentation for its
tagger).
\item \textbf{Buffer globally.}  Tags are dilated by
\code{TagBuffer} on the \emph{level-wide} mask, so a buffer that
overflows a patch edge lands in the neighboring patch --- the Fortran
\code{BufferAndList}/\code{InflateEdgeTags} broadcast, including the
periodic wrap where applicable.
\item \textbf{Clip for proper nesting} (\cref{sec:nesting}).
\item \textbf{Cluster per parent.}  Berger--Rigoutsos signature
splitting runs \emph{independently on each patch's own tags}
(\code{RefineGrid} passes \code{NewSubGrids} the single patch's
\code{errorflags}); candidate rectangles start from the tags' bounding
box inside the patch and only ever shrink or split, so no cluster box
can leave its parent --- \cref{inv:parent} holds by construction.
\item \textbf{Refine.}  Each accepted box becomes one child (its exact
$2\times$ refinement), tagged with its parent's index.
\end{enumerate}

Only the buffering is global; the clustering is local.  As in both
parent codes there is \emph{no aspect-ratio constraint} --- boxes are
free rectangles subject to \code{MinBoxWidth}, the fill-ratio target,
and the nesting clip.

\begin{cautionbox}
Per-parent clustering fragments features that cross patch boundaries.
Occasionally a fragment is a sliver (one cell wide after the nesting
clip) that cannot grow within its parent; such boxes are merged into an
adjacent box of the same parent when possible and otherwise dropped
with a warning (\code{mergeOrDropThin}) --- locally \emph{less}
refinement, never a fragile patch.  The append-only elliptic driver
does not revisit dropped cells on older levels; coverage must be
checked when resolving a prescribed region is required.  The nesting
clip keeps the mesh 1-conforming.  The Fortran instead grows such boxes over untagged
parent cells (its trim floor \code{minimumgridpoints} without an
allowed-region concept), refining slightly more in these spots.
\end{cautionbox}

\section{Proper nesting without restarts}
\label{sec:nesting}

The Fortran enforces 2-level conformity \emph{a posteriori}: a
patch-resolution level map (\code{levellandscape}) detects level-$L$
patches adjacent to level-$(L{-}2)$ regions, and the mesh build restarts
from level $L{-}2$ with augmented tags (\code{FindMeshDefects}, up to
100 \code{amrrestarts}).  The MATLAB tree solvers instead \emph{prevent}
defects: tags are clipped to cells whose full 8-neighborhood is covered
by the current finest level (\code{erode8}; x-periodic in the
CH--Stokes solver).  Consequences, all verified by the tree checkers:

\begin{itemize}[leftmargin=2em]
\item every ring cell of a child is covered by the parent level's
patch union (or lies beyond a physical boundary), so the QIF stencil
always reads synchronized coarse data;
\item interface jumps are one level (1-conforming composite mesh);
\item no rebuild loops --- each regrid pass is a single sweep.
\end{itemize}

% ======================================================================
\chapter{Single-Parent Transfers in the Elliptic Solver}
\label{ch:ebt}

\code{EllipticBSAMTree} keeps the discretization, kernels, smoother,
QIF/LIF interpolation, flux-balance corrections, composite residual and
direct coarsest solve of \code{EllipticBSAM} \emph{unchanged}.  This
chapter records only the moving parts that were rebuilt on the tree.

\section{What changed relative to \code{EllipticBSAM}}

\begin{center}\small
\begin{tabular}{@{}p{0.3\textwidth}p{0.31\textwidth}p{0.31\textwidth}@{}}
\toprule
 & \code{EllipticBSAM} & \code{EllipticBSAMTree} \\
\midrule
clustering & level-global B--R & per-parent B--R (\cref{sec:perparent}) \\
lineage & \code{parents\{p\}} list & \code{parent(p)} scalar + \code{mbounds} \\
window/ring & per-cell gathers from all owners & one padded-parent slice \\
restriction, FAS load & scatters over all intersecting owners & single-rect scatter into the parent \\
coefficient fill-down & \code{rectIsect} over all coarse patches & sliced write into the parent + shared-face push (\cref{sec:sharedface}) \\
correction targets & owner-map scatter & unchanged (targets may be the parent's siblings) \\
\bottomrule
\end{tabular}
\end{center}

The flux-balance correction targets are the one place the tree still
routes through an owner map: when a child touches its parent's edge,
the coarse cell just outside the footprint belongs to a \emph{sibling
of the parent}.  This is precomputed (\code{ops.corrW/E/S/N}), the
static analogue of the Fortran's global \code{masscorrlayer} strip
list.

\section{The duplicated shared-face lesson}
\label{sec:sharedface}

\begin{cautionbox}
Coarse patch boundary \emph{faces} are stored twice --- once in each
abutting patch's face-coefficient array.  When a child's footprint
touches its parent's edge, the coefficient fill-down (coarse face $D$
$\leftarrow$ mean of the two overlying fine faces) must write the
averaged value into \emph{both} copies, or the two sides of that face
compute different fluxes and the discrete conservation identity breaks.
\code{ebtCoeffFilldown} pushes boundary-face writes to the abutting
siblings through the owner map (\code{pushFaceX/pushFaceY}).  The
symptom when this is missing is instructive: a mass defect that
\emph{floors} (does not shrink with tighter solver tolerance) and
appears only for variable $D$ on adaptive meshes --- constant
coefficients hide it, because the fill-down average then equals the
analytic value.
\end{cautionbox}

\section{Verification}
\label{sec:ebtverify}

\code{tests/runAllTreeTests} reruns the entire \code{EllipticBSAM}
suite (T1--T8) on the tree solver and adds T9:

\begin{itemize}[leftmargin=2em]
\item \textbf{T1--T7 parity.}  Identical numbers to the QC-ring solver:
second-order $L^\infty$ rates (1.98--2.03) on uniform, box-in-box,
L-shaped and fully adaptive hierarchies for QIF and LIF;
$h$- and depth-independent V(3,3) factors $0.026$--$0.054$;
conservation defects $7\cdot 10^{-17}$--$3\cdot 10^{-15}$.
\item \textbf{T8 robustness.}  24 random tag disks; per-parent
clustering yields 50 patches (QC-ring: 49) of minimum width 4, factor
$0.052$, mass defect $1.1\cdot 10^{-13}$.
\item \textbf{T9 tree invariants.}  \code{ebtTreeCheck} passes on the
adaptive hierarchy, and the engineered straddling scenario (a tag
centered on the shared edge of two coarse patches) produces exactly two
children with parents $[1\ 2]$ --- never a multi-parent child --- with
V-cycle factor $0.033$ and mass defect $9.5\cdot 10^{-13}$.
\end{itemize}

\begin{remark}
The companion audit measured identical per-cycle factors for a
multi-parent child (QC ring) and a single-parent child on the same
mesh.  The tree buys structure and simplicity, not V-cycle speed; its
performance value lies in the O(1) sliced transfers and in what it
enables next --- the CH--Stokes solver's face-centered machinery would
be substantially harder against multi-parent children.
\end{remark}

% ======================================================================
\chapter{Verification}
\label{ch:verify}

\section{EllipticBSAMTree: runAllTreeTests}

T1--T8 reproduce the \code{EllipticBSAM} suite on the tree
(\cref{sec:ebtverify}); T9 validates the tree invariants and the
per-parent splitting of a straddling feature.  Headline numbers
(MATLAB R2026a): uniform and adaptive $L^\infty$ rates 1.98--2.03;
V(3,3) factors 0.026--0.054 independent of $h$ and depth; conservation
defects at roundoff ($10^{-16}$--$10^{-13}$); suite time
$\approx37$\,s.

\chapter{Mathematical and Diagnostic Clarifications}
\label{ch:audit2026}
\section{Compatibility is a source projection}
For all-Neumann data and $C\equiv0$, the discrete source and the outward
boundary derivative must satisfy
\[
 \sum_{i\in\mathcal V} h_{x,i}h_{y,i} f_i
 +\sum_{F\subset\partial\Omega} A_F D_F g_F=0.
\]
\code{ebtCompatProject} subtracts the defect divided by the visible
domain area from the original sampled source on every level.  The
removed constant is \code{H.compatibility.sourceShift}; inspect it
because incompatible data are thereby changed.  Per-level
\code{sourceShift} bookkeeping ensures that previously projected
coarse data and fresh fine data are treated consistently at every
refinement.  Higher source quadrature improves the source integral;
it is not exact for arbitrary data and does not by itself enforce the
boundary-flux compatibility condition.

\section{Reaction fixes the mean}
Since $\mathcal L(u+a)=\mathcal Lu+Ca$, only the all-Neumann,
$C\equiv0$ problem has an additive-constant gauge.  The error routine
aligns the error mean exactly in this singular case, for every
\code{SourceQuad}; otherwise it retains the absolute error.  The
removed shift is \code{err.meanShift}.  A small positive reaction is
not classified as zero using an absolute numerical threshold.
\code{testEbtNeumannSemantics} verifies this rule and repeated
refinement of a compatible quartic manufactured solution.

\section{Proof boundaries for the actual QIF stencil}
The tree changes data movement, not the stencil coefficients.  At a
regular west QIF face,
\[
 Q_{1,3}=\tfrac8{15}Q^c_{1,2}+\tfrac23Q_{2,4}-\tfrac15Q_{3,5}.
\]
The equation for $Q_{2,3}$ therefore contains
$+D Q_{3,5}/(5h_x^2)$.  This is a positive off-diagonal coefficient
even for constant $D$ on a square grid, since $Q_{3,5}$ is not an
ordinary nearest neighbor of $Q_{2,3}$.  Blanket $Z$-matrix and
$M$-matrix claims for the fully eliminated multidimensional QIF
operator are consequently unjustified.  Conservation follows from
matching face fluxes; a comparison principle, energy estimate and
uniform V-cycle bound need separate arguments.  The companion theory
document's multidimensional sign lemma and conclusions depending on
it remain proof-review items.

The physical boundary fold gives exact point updates with neighbors
fixed.  At coarse--fine interfaces, ghost interpolation is lagged
within each color pass; global parity alone does not make this literal
Gauss--Seidel on the fully eliminated, wider QIF stencil.  The reported
contraction factors establish performance on the tested mesh families.

\section{RL environment: what is measured}
The \code{rl/} API can observe an unsolved root
(\code{SolveRoot=false}), take individual cycles, and refine or polish
the current hierarchy.  A stage observation supersedes the residual
of any preceding single cycle.  \code{testRlState} verifies these
transitions and rejection of non-finite or fractional box indices.

The reward's relative-truncation-error proxy omits a two-cell margin
around every patch and all patches narrower than eight cells.  A
smaller proxy can therefore partly reflect a changed measurement
region.  It is not a guaranteed discretization-error estimator; report
true manufactured error, residual, visible cells and work alongside it.
For noisy numeric forcing, comparison with the clean generating field
includes data/interpolation error and is not error against an exact
solution of the perturbed PDE.  Calibration-fit correlations describe
the sampled calibration set, not an out-of-sample error guarantee.

\section{Headless regression command}
\begin{lstlisting}[language=bash]
/Applications/MATLAB_R2026a.app/bin/matlab -batch "maxNumCompThreads(2); cd('EllipticBSAMTree/tests'); unitTreeKernels; runAllTreeTests; stressTreeInterfaces; testRlApi; testRlState; testEbtNeumannSemantics"
\end{lstlisting}
Keep the tree tests on an isolated MATLAB path to avoid the other
package's \code{mmsCases}.  Historical positive-reaction Neumann error
tables may use a mean-aligned error; rerun before comparing with the
corrected absolute norms.  Dated audit logs preserve both versions.

\chapter{Function Reference}

\section{EllipticBSAMTree (\code{ebt*})}

\begin{center}\small
\begin{tabularx}{\textwidth}{@{}>{\raggedright\arraybackslash}p{0.50\textwidth}X@{}}
\toprule
function & responsibility \\
\midrule
\code{ebtSolve} & driver: setup, solve--tag--refine loop, outputs \\
\code{ebtProblem}, \code{ebtOptions} & problem and option structs \\
\code{ebtSetup}, \code{ebtMakeLevel} & hierarchy and per-patch fields (+ lineage) \\
\code{ebtAddLevel} & append a level from per-parent cluster boxes \\
\code{ebtTagCluster} & tagging, global buffer, nesting clip, per-parent B--R \\
\code{ebtBuildOps} & owner maps, same-level exchange, single-parent descriptors \\
\code{ebtFillGhosts}, \code{ebtInterpCf} & ghost priority order; QIF/LIF stencils \\
\code{ebtTransfer} & restrict / ring / correct / filldown \\
\code{ebtVcycle}, \code{ebtSolveCycles}, \code{ebtRelaxLevel} & AFAS engine \\
\code{ebtCoarseLoad}, \code{ebtFaceCorrections} & FAS load + flux-balance corrections \\
\code{ebtCoeffFilldown}, \code{ebtParentCoeffs}, \code{ebtGatherParentCoeffs} & coefficient compatibility \\
\code{ebtCoarsest}, \code{ebtKernels}, \code{ebtBcFold} & level-1 solve; kernels; boundary fold \\
\code{ebtCompResidual}, \code{ebtErrorNorms}, \code{ebtMassCheck} & diagnostics \\
\code{ebtTreeCheck} & tree invariant validator \\
\code{ebtPlot}, \code{ebtSample}, \code{ebtRte}, \code{ebtCompositeMean}, \code{ebtGatherWindow} & utilities \\
\bottomrule
\end{tabularx}
\end{center}

\begin{thebibliography}{9}
\bibitem{feng2018}
W.~Feng, Z.~Guo, J.~S.~Lowengrub, S.~M.~Wise,
\emph{A mass-conservative adaptive FAS multigrid solver for
cell-centered finite difference methods on block-structured,
locally-cartesian grids},
J.~Comput.~Phys.\ \textbf{352} (2018) 463--497.
\bibitem{br1991}
M.~J.~Berger, I.~Rigoutsos,
\emph{An algorithm for point clustering and grid generation},
IEEE Trans.\ Systems, Man, and Cybernetics \textbf{21} (1991)
1278--1286.
\bibitem{vanka1986}
S.~P.~Vanka,
\emph{Block-implicit multigrid solution of Navier-Stokes equations in
primitive variables},
J.~Comput.~Phys.\ \textbf{65} (1986) 138--158.
\end{thebibliography}

\end{document}
